\documentclass[10pt,a4paper]{amsart}
\usepackage[margin=19mm]{geometry}
\usepackage{amsmath,amssymb,mathtools}
\usepackage[scr=rsfs,cal=euler]{mathalfa}
\usepackage{xfrac}
\usepackage[unicode,hidelinks,bookmarks=false]{hyperref}
\newcommand{\ie}{i.e.\ }
\newcommand{\cf}{cf.\ }
\newcommand{\resp}{respectively\ }
\newcommand{\Subsection}{\subsection}
\providecommand{\nobreakdash}{\nobreak\hbox{-}}
\setlength{\emergencystretch}{2em}
\multlinegap0pt
\usepackage{bm}
\usepackage{bbm}
\usepackage{dsfont}
\usepackage{xspace}
\usepackage{cancel}
\usepackage{mathtools}
\usepackage{amsthm}
\makeatletter
\@ifundefined{theorem}{\newtheorem{theorem}{Theorem}}{}
\@ifundefined{lemma}{\newtheorem{lemma}[theorem]{Lemma}}{}
\@ifundefined{remark}{\newtheorem{remark}[theorem]{Remark}}{}
\makeatother
\DeclareMathOperator{\Prec}{B}
\DeclareMathOperator{\Succ}{A}
\DeclareMathOperator{\str}{str}
\title[A note on the strength of a hypercube]{A note on the strength of a hypercube}
\author{Melissa A. Huggan, M.E. Messinger, Dylan Pearson}
\date{Source: arXiv:2507.21908v1 (CC BY 4.0)}
\begin{document}
\maketitle

\noindent\emph{Source and licence.} A note on the strength of a hypercube. Version arXiv:2507.21908v1 (CC BY 4.0), distributed under the Creative Commons Attribution 4.0 International licence, as recorded in the arXiv licence metadata for this exact version and shown on the abstract page https://arxiv.org/abs/2507.21908v1. Full rights notice: licenses/arxiv-2507.21908-CC-BY-4.0.txt.\par\smallskip

\noindent\textit{Editorial scope.} Counted unit: the Theorem whose source label is \texttt{thm:both10even} in the article (source0.tex lines 370--370, source label \texttt{thm:both10even}), delivered together with the article's own complete original proof of it (source0.tex lines 373--416, one \texttt{proof} environment of 44 lines). No mathematics is written, repaired or re-derived: statement and proof are the article's own text line for line. Required context retained uncounted: remark1 (lines 252--263); lem:ends (lines 266--267); obs:x'y' (lines 290--290); lemma0.1 (lines 359--362). Every definition, display and label of those lines is retained unchanged. Omitted from this package and not redistributed: 1--251, 264--265, 268--289, 291--358, 363--369, 371--372, 417--492. Nothing inside a retained excerpt is omitted or altered: every retained body is delivered exactly as the article's lines are written, so no omission receipt of any kind accompanies this package. Citations: the retained excerpts carry no citation command at all, so no bibliography entry of the article is transcribed and no reference is redistributed. Reference adaptation: none was needed and none was made; every reference inside the delivered unit resolves inside it, so no unresolved reference or citation is produced. Printed slips kept literal: no printed word, symbol, hypothesis or number of the counted statement or of its proof was altered. Layout and class: the article's own class is \texttt{article} with packages including amsmath, amssymb, xcolor, url, amsthm, enumerate, caption, subcaption, soul, geometry; the wrapper is the lane's pinned \texttt{amsart} editorial wrapper at 10pt on A4, which restates the theorem-like environments the retained text needs and adds the article's own macro definitions that the retained text needs (\texttt{\textbackslash Prec}, \texttt{\textbackslash Succ}, \texttt{\textbackslash str}), with extra wrapper packages (bm, bbm, dsfont, xspace, cancel, mathtools). The delivered numbering is the unit's own. Assets: no retained span contains a figure environment, a graphics-inclusion command, a PDF-page inclusion, a table or a TikZ picture; no image file is copied into this package, the package declares no asset dependency and no image file is redistributed with it. Licence: version arXiv:2507.21908v1 of this article is distributed under the Creative Commons Attribution 4.0 International licence, as recorded in the arXiv licence metadata for this exact version and shown on the abstract page https://arxiv.org/abs/2507.21908v1. A full rights notice is delivered at licenses/arxiv-2507.21908-CC-BY-4.0.txt. Characters: every retained excerpt is pure ASCII, so the delivered engine, XeLaTeX with its default Latin Modern text font, reports no missing character.
\begin{remark}\label{remark1}
For any integer $n \geq 1$, let $x$ and $y$ be two $n$-bit strings. The vertices of $Q_n$ that correspond to $x$ and $y$ are adjacent in $Q_n$ if and only if $x$ and $y$ have Hamming distance one. Fixing a bijective function $f$ from the $n$-bit strings to $\{1, 2, \ldots, 2^n\}$ necessarily implies that there exist $n$-bit strings $x$ and $y$ such that $\str_f(Q_n) = f(x)+f(y)$. 

For a pair $x,y$ of $n$-bit strings for which  $\str_f(Q_n)=f(x)+f(y)$ and $x$ has weight $i$, we make the following assumption throughout this note: 

\begin{center}
    $x$ has odd weight $i$ and $y$ has even weight $i-1$.
\end{center}

\noindent
Since $x$ and $y$ have Hamming distance one, their weights will be of different parity; thus we may assume $i$ is odd.  Consequently, $y$ must have even weight $i-1$ or $i+1$.  Strings of weight $i-1$ fall later in sequence $S_n$ (and therefore map to larger function values) than strings of weight $i+1$. Thus, if $x$ has odd weight $i$ then $y$ has even weight $i-1$.
 \end{remark}

\begin{lemma}\label{lem:ends} For any integer $n \geq 2$, there exists an $(n-1)$-bit string $w$ such that $x=w1$, $y=w0$ and $\str_f(Q_n) = f(x)+f(y)$.
\end{lemma}

\begin{lemma}\label{obs:x'y'} For any integer $n \geq 2$, let $w$ be an $(n-1)$-bit string of weight $i-1$ for some odd integer $i$.  Then $$f(w1)+f(w0) = 2^n + \tbinom{n-1}{i} - |\Succ(w1,R_n^i)|+|\Prec(w0,S_n^{i-1})|+1.$$\end{lemma}

\begin{lemma}\label{lemma0.1} Let $n-2$ be an even positive integer and $j: 0\leq j \leq n-2$, with $j \neq \frac{n-2}{2}$ and $S_{n-2}^j = (v_1,v_2,\dots,v_{k-1},v_k)$, where $k = \binom{n-2}{j}$. Then \begin{itemize}
\item $S_{n-2}^{n-j-2} = (\overline{v_k},\overline{v_{k-1}},\dots,\overline{v_2},\overline{v_1})$;
\item $R_{n-2}^j = (v_k,v_{k-1},\dots,v_2,v_1)$; and 
\item $R_{n-2}^{n-j-2}=(\overline{v_1},\overline{v_2},\dots,\overline{v_{k-1}},\overline{v_k}).$\end{itemize}\end{lemma}

\begin{theorem}\label{thm:both10even} For any integer $n \geq 5$, let $x$ and $y$ be $n$-bit strings for which $\str_f(Q_{n}) = f(x)+f(y)$.  If $x$ and $y$ both begin with $10$ and $n$ is even, then $$\str_f(Q_{n}) \leq \str_f(Q_{n-2}) + 3 \cdot 2^{n-2} + \binom{n-3}{\lceil \frac{n-3}{2}\rceil}+\binom{n-2}{\lceil \frac{n-2}{2}\rceil}.$$ \end{theorem} 

\begin{proof} Let $x$ and $y$ be $n$-bit strings for which $\str_f(Q_{n}) = f(x)+f(y)$, and assume $n$ is even, and $x$ and $y$ both begin with $10$.  By Remark~\ref{remark1}, $x$ has weight $i$ and $y$ has weight $i-1$.  By Lemma~\ref{lem:ends}, $x$ and $y$ differ only in that the last bit in $x$ is $1$ and the last bit in $y$ is $0$.  Thus, $x = 10w1$ and $y=10w0$ for some string $w$ of length $n-3$ and weight $i-2$.  We will bound the difference $(f(x)+f(y)) - (f(\overline{w}1)+f(\overline{w}0))$ because we can exploit the relationships between $f(x)$ and $f(\overline{w}0)$ and between $f(y)$ and $f(\overline{w}1)$.  

By Lemma~\ref{obs:x'y'}, $$f(x)+f(y) = 2^{n}+\tbinom{n-1}{i}-|\Succ(x,R_{n}^i)|+|\Prec(y,S_{n}^{i-1})|+1$$ and since $\overline{w}$ is an $(n-3)$-bit string of weight $n-i-1$, by the same lemma: $$f(\overline{w}1)+f(\overline{w}0)=2^{n-2} + \tbinom{n-3}{n-i}-|\Succ(\overline{w}1,R_{n-2}^{n-i})|+|\Prec(\overline{w}0,S_{n-2}^{n-i-1})|+1.$$  Thus,\medskip 

\noindent $\Big(f(x)+f(y)\Big) - \Big(f(\overline{w}1)+f(\overline{w}0)\Big)$ $$=3 \cdot 2^{n-2}+\tbinom{n-1}{i}-\tbinom{n-3}{n-i} -|\Succ(x,R_{n}^i)|+|\Prec(y,S_{n}^{i-1})| +|\Succ(\overline{w}1,R_{n-2}^{n-i})|-|\Prec(\overline{w}0,S_{n-2}^{n-i-1})|$$ \begin{equation} =3 \cdot 2^{n-2} + \tbinom{n-1}{i}-\tbinom{n-3}{i-3}
-|\Succ(x,R_{n}^i)|+|\Prec(y,S_{n}^{i-1})|+|\Succ(\overline{w}1,R_{n-2}^{n-i})|-|\Prec(\overline{w}0,S_{n-2}^{n-i-1})|.
\label{eqtogether} \end{equation}

\bigskip

Our goal is to simplify Equation~(\ref{eqtogether}).  We first observe that \begin{equation}\label{eq2nd} |\Prec(\overline{w}0,S_{n-2}^{n-i-1})|+|\Succ(\overline{w}0,S_{n-2}^{n-i-1})|+1 = \tbinom{n-2}{n-i-1} = \tbinom{n-2}{i-1}.\end{equation}
We next relate $|\Succ(x,R_{n}^i)|$ to $|\Prec(\overline{w}0,S_{n-2}^{n-i-1})|$.  The sequence $R_{n}^i$ consists of \begin{itemize}
\item sequence $R_{n-2}^{i-2}$ with $11$ prefixed to each string in the sequence; followed by\vspace{-0.1in}
\item sequence $R_{n-2}^{i-1}$ with $10$ prefixed to each string in the sequence; followed by\vspace{-0.1in}
\item sequence $R_{n-2}^{i-1}$ with $01$ prefixed to each string in the sequence; followed by\vspace{-0.1in}
\item sequence $R_{n-2}^i$ with $00$ prefixed to each string in the sequence.
\end{itemize}
Then \begin{eqnarray} |\Succ(x,R_{n}^i)| &=& |\Succ(10w1,R_{n}^i)| \nonumber \\ &=& |\Succ(w1,R_{n-2}^{i-1})|+\tbinom{n-2}{i-1}+\tbinom{n-2}{i} \nonumber\\ 
&=& |\Succ(\overline{w}0,S_{n-2}^{n-i-1})|+\tbinom{n-2}{i-1}+\tbinom{n-2}{i} ~~~\text{by Lemma~\ref{lemma0.1}}\nonumber \\ &=& 2\tbinom{n-2}{i-1}+\tbinom{n-2}{i}-1-|\Prec(\overline{w}0,S_{n-2}^{n-i-1})| ~~~\text{by Equation~(\ref{eq2nd}).}\label{eq1st}\end{eqnarray}  Equation~(\ref{eq1st}) establishes the relationship between $|\Succ(x,R_{n}^i)|$ and $|\Prec(\overline{w}0,S_{n-2}^{n-i-1})|$.  



We now observe that \begin{equation} |\Succ(\overline{w}1,R_{n-2}^{n-i})|+|\Prec(\overline{w}1,R_{n-2}^{n-i})|+1 = \tbinom{n-2}{n-i} = \tbinom{n-2}{i-2}\label{eqzz} \end{equation} and we next relate $|\Prec(y,S_{n}^{i-1})|$ to $|\Succ(\overline{w}1,R_{n-2}^{n-i})|$. The sequence $S_{n}^{i-1}$ consists of \begin{itemize}
\item sequence $S_{n-2}^{i-1}$ with $00$ prefixed to each string in the sequence; followed by\vspace{-0.1in}
\item sequence $S_{n-2}^{i-2}$ with $01$ prefixed to each string in the sequence; followed by\vspace{-0.1in}
\item sequence $S_{n-2}^{i-2}$ with $10$ prefixed to each string in the sequence; followed by\vspace{-0.1in}
\item sequence $S_{n-2}^{i-3}$ with $11$ prefixed to each string in the sequence.
\end{itemize}

Then \begin{eqnarray} |\Prec(y,S_{n}^{i-1})| &=& |\Prec(10w0,S_{n}^{i-1})| \nonumber \\  
&=& |\Prec(w0,S_{n-2}^{i-2})|+\tbinom{n-2}{i-2}+\tbinom{n-2}{i-1} \nonumber \\
&=& |\Prec(\overline{w}1,R_{n-2}^{n-i})|+\tbinom{n-2}{i-2}+\tbinom{n-2}{i-1} ~~~\text{by Lemma~\ref{lemma0.1}}\nonumber \\
&=& \tbinom{n-2}{i-2}-|\Succ(\overline{w}1,R_{n-2}^{n-i})|-1+\tbinom{n-2}{i-2}+\tbinom{n-2}{i-1} ~~~\text{by Equation~(\ref{eqzz}) \nonumber}\\
&=& 2\tbinom{n-2}{i-2}+\tbinom{n-2}{i-1}-1-|\Succ(\overline{w}1,R_{n-2}^{n-i})|.\label{eqczz}
\end{eqnarray}



We use Equations~(\ref{eq1st}) and~(\ref{eqczz}) to simplify the expression in Equation~(\ref{eqtogether}) to:
\begin{eqnarray*} \Big(f(x)+f(y)\Big) - \Big(f(\overline{w}1)+f(\overline{w}0)\Big) &=& 3 \cdot 2^{n-2} + \tbinom{n-1}{i}-\tbinom{n-3}{i-3}+2\tbinom{n-2}{i-2}-\tbinom{n-2}{i-1}-\tbinom{n-2}{i}\\ 
&=& 3 \cdot 2^{n-2} + 2\tbinom{n-2}{i-2}-\tbinom{n-3}{i-3} \\ &=& 3 \cdot 2^{n-2} +\tbinom{n-2}{i-2}+\tbinom{n-3}{i-2} \\ &\leq& 3 \cdot 2^{n-2} + \tbinom{n-2}{\lceil (n-2)/2\rceil }+\tbinom{n-3}{\lceil (n-3)/2\rceil}.\end{eqnarray*}

So $\str_f(Q_{n}) \leq \str_f(Q_{n-2}) + 3 \cdot 2^{n-2} +\tbinom{n-2}{\lceil (n-2)/2\rceil }+\tbinom{n-3}{\lceil (n-3)/2\rceil}$.
\end{proof}

\end{document}
